% 第 14 章 移动算法和迭代器（由 tools/convert.py 从原 content/ 自动生成，勿手改）

C++标准库对迭代元素时移动元素提供了特殊支持。为此，C++11中引入了特殊算法和特殊迭代器(移动迭代器)。本章讨论如何使用它们。\par


\section{移动算法}

C++标准库提供了一些移动元素的算法。这些算法是:\par

\begin{itemize}
  \item \textbf{std::move()}，将元素移动到另一个范围或在同一范围内向后移动(不要将此算法与\textit{std::move()}标记不再需要其值的对象混淆)\par
  指定目标范围的开头，元素将从源范围的开头移动到末尾。
  \item \textbf{std::\allowbreak{}move\_\allowbreak{}backward(\allowbreak{})\allowbreak{}}，将元素移动到另一个范围或在同一范围内向前移动\par
  指定目标范围的结束，元素将从源范围的结束移动到开始。
\end{itemize}

这些算法使用的是\textit{std::copy()}和\textit{std::\allowbreak{}copy\_\allowbreak{}backward(\allowbreak{})\allowbreak{}}算法的对等操作。是的，\textit{std::move()}还有另一个重载，接受多个形参(三个迭代器(从C++17起)，还有一个可选的并行执行策略)。\par

这些算法的效果是在迭代每个元素时调用\textit{std::\allowbreak{}move(\allowbreak{}elem)\allowbreak{}}对目标范围进行移动赋值。\par

考虑以下例子:\par

\begin{codeboxt}[cpp]{lib/movealgo.cpp}
#include <iostream>
#include <string>
#include <vector>
#include <list>
#include <algorithm>

template<typename T>
void print(const std::string& name, const T& coll)
{
  std::cout << name << " (" << coll.size() << " elems): ";
  for (const auto& elem : coll) {
    std::cout << " '" << elem << "'";
  }
  std::cout << "\n";
}

int main(int argc, char** argv)
{
  std::list<std::string> coll1 { "love", "is", "all", "you", "need" };
  std::vector<std::string> coll2;

  // ensure coll2 has enough elements to overwrite their values:
  coll2.resize(coll1.size());

  // print out size and values:
  print("coll1", coll1);
  print("coll2", coll2);
  // move assign the values from coll1 to coll2
  // - not changing any size
  std::move(coll1.begin(), coll1.end(), // source range
  coll2.begin()); // destination range

  // print out size and values:
  print("coll1", coll1);
  print("coll2", coll2);

  // move assign the first three values inside coll2 to the end
  // - not changing any size
  std::move_backward(coll2.begin(), coll2.begin()+3, // source range
  coll2.end()); // destination range

  // print out size and values:
  print("coll1", coll1);
  print("coll2", coll2);
}
\end{codeboxt}

当调用\textit{std::move()}时，将源容器的所有值赋给目标容器:\par

\begin{codebox}[cpp]
// move assign the values from coll1 to coll2
// - not changing any size
std::move(coll1.begin(), coll1.end(), // source range
      coll2.begin()); // destination range
\end{codebox}

与覆盖算法一样，目标容器必须有足够的元素(否则就会有未定义的行为)。元素的数量不会改变(无论是在源范围还是在目标范围)。但是，源范围的元素处于“已移动”状态。因此，这个调用之后，就不知道源范围内字符串的值(除非为已移动对象指定了行为，比如只移动类型)。\par

当调用\textit{std::\allowbreak{}move\_\allowbreak{}backward(\allowbreak{})\allowbreak{}}算法时，将前三个元素赋值到同一个集合的末尾:\par

\begin{codebox}[cpp]
// move assign the first three values inside coll2 to the end
// - not changing any size
std::move_backward(coll2.begin(), coll2.begin()+3, // source range
           coll2.end()); // destination range
\end{codebox}

同样，如果没有另一个值移走，元素的状态将会是已移动。因此，不再知道前两个元素的值(第三个元素的值被第一个元素的移动赋值覆盖)。在这通调用之后，我们只知道\textit{coll2}的最后三个元素是\textit{love}, \textit{is}和\textit{all}。\par

因此，整个程序的输出是这样的(?表示不确定的值):\par

\begin{myOutput}
  coll1 (5 elems): 'love' 'is' 'all' 'you' 'need' \\
  coll2 (5 elems): '' '' '' '' '' \\
  coll1 (5 elems): '?' '?' '?' '?' '?' \\
  coll2 (5 elems): 'love' 'is' 'all' 'you' 'need' \\
  coll1 (5 elems): '?' '?' '?' '?' '?' \\
  coll2 (5 elems): '?' '?' 'love' 'is' 'all'
\end{myOutput}

字符串移动后通常是空的，但这不能保证。实践中，我甚至在某个平台上发现了这样的输出:\par

\begin{myOutput}
  coll1 (5 elems): 'love' 'is' 'all' 'you' 'need' \\
  coll2 (5 elems): '' '' '' '' '' \\
  coll1 (5 elems): '' '' '' '' '' \\
  coll2 (5 elems): 'love' 'is' 'all' 'you' 'need' \\
  coll1 (5 elems): '' '' '' '' '' \\
  coll2 (5 elems): 'need' 'you' 'love' 'is' 'all'
\end{myOutput}

所以，与往常一样:使用已移动元素时要谨慎。\par


\section{移除性算法}

根据设计，C++算法使用迭代器来处理容器和范围的元素。然而，就像指针操作数组一样，迭代器只能读取和写入值，不能插入或删除元素。因此，“删除”算法并不是真正删除元素，只是将所有未移除元素的值移动到已处理范围的前面。\par

例如，给定以下整数序列:\par

\begin{myOutput}
1 2 3 4 5 4 3 2 1
\end{myOutput}

调用\textit{std::\allowbreak{}remove(\allowbreak{})\allowbreak{}}算法来删除值为2的所有元素，修改后的序列:\par

\begin{myOutput}
1 3 4 5 4 3 1 2 1
\end{myOutput}

所有不是2的元素都移到前面，作为新的结束(“last”元素后面的位置)，返回2的位置。\par

可能的话，这些算法会移动，会以已移状态保留元素。本例中，如果元素是字符串，那么最后2个元素将保持不变，但最后1个元素将向前移动到2之前，以便最后一个元素处于已移动状态。\par

因此，这些算法也可以将元素保留为已移动状态。事实上，以下算法可以创建已移动状态:\par

\begin{itemize}
  \item \textit{std::\allowbreak{}remove(\allowbreak{})\allowbreak{}}和\textit{std::\allowbreak{}remove\_\allowbreak{}if(\allowbreak{})\allowbreak{}}
  \item \textit{std::\allowbreak{}unique(\allowbreak{})\allowbreak{}}
\end{itemize}

看一个完整的例子，可以看到元素是否为已移动状态:\par

\begin{codeboxt}[cpp]{lib/email.hpp}
#include <iostream>
#include <cassert>
#include <string>

// class for email addresses
// - asserts that each email address has a @
// - except when in a moved-from state
class Email {
  private:
  std::string value; // email address
  bool movedFrom{false}; // special moved-from state
  public:
  Email(const std::string& val)
  : value{val} {
    assert(value.find('@') != std::string::npos);
  }
  Email(const char* val) // enable implicit conversions for string literals
  : Email{std::string(val)} {
  }
  std::string getValue() const {
    assert(!movedFrom); // or throw
    return value;
  }
  ...
  // implement move operations to signal a moved-from state:
  Email(Email&& e) noexcept
  : value{std::move(e.value)}, movedFrom{e.movedFrom} {
    e.movedFrom = true;
  }
  Email& operator=(Email&& e) noexcept {
    value = std::move(e.value);
    movedFrom = e.movedFrom;
    e.movedFrom = true;
    return *this;
  }
  // enable copying:
  Email(const Email&) = default;
  Email& operator=(const Email&) = default;

  // print out the current state (even if it is a moved-from state):
  friend std::ostream& operator<< (std::ostream& strm, const Email& e) {
    return strm << (e.movedFrom ? "MOVED-FROM" : e.value);
  }
};
\end{codeboxt}

通过实现移动构造函数和移动赋值操作符，设置了一个成员\textit{movedFrom}，并在输出操作符中求值。\par

现在让我们将该类中的一些元素放入vector中，并使用算法\textit{std::\allowbreak{}remove\_\allowbreak{}if(\allowbreak{})\allowbreak{}}删除一些元素。本例中，我们删除所有以".de"结尾的电子邮件地址:\par

\begin{codeboxt}[cpp]{lib/removeif.cpp}
#include "email.hpp"
#include <iostream>
#include <string>
#include <vector>
#include <algorithm>
int main()
{
  std::vector<Email> coll{ "tomdomain.de", "jillcompany.com",
    "sarahdomain.de", "hanacompany.com" };

  // remove all email addresses ending with ".de":
  auto newEnd = std::remove_if(coll.begin(), coll.end(),
  [] (const Email& e) {
    auto&& val = e.getValue();
    return val.size() > 2 &&
    val.substr(val.size()-3) == ".de";
  });

  // print elements up to the new end:
  std::cout << "remaining elements:\n";
  for (auto pos = coll.begin(); pos != newEnd; ++pos) {
    std::cout << " '" << *pos << "'\n";
  }

  // print all elements in the container:
  std::cout << "all elements:\n";
  for (const auto& elem : coll) {
    std::cout << " '" << elem << "'\n";
  }
}
\end{codeboxt}

程序输出如下:\par

\begin{myOutput}
remaining elements: \\
"jill@company.com" \\
"hana@company.com" \\
all elements: \\
"jill@company.com" \\
"hana@company.com" \\
"sarah@domain.de" \\
"MOVED-FROM"
\end{myOutput}

“删除”以“.de”结尾的元素之后，新的结束位置是第二个元素后面的位置。但在容器中，第三个元素没有移动，第四个元素处于已移动状态，因为它移动赋值给了第二个元素(之前移动到第一个元素)。\par

使用这些已移动对象时要小心，原因在关于已移动状态的章节中已经讨论过了。\par

\mySmallsection{14.3.1 移动算法中的迭代器}

算法中使用移动迭代器，通常只有在算法保证每个元素只使用一次时才有意义。因此，该算法应该:\par

\begin{itemize}
  \item 要求源的输入迭代器类别和目标的输出迭代器类别
  \item 或者保证每个元素只使用一次(例如，为\textit{std::\allowbreak{}for\_\allowbreak{}each(\allowbreak{})\allowbreak{}}算法指定)
\end{itemize}

对于具有可调用对象的算法，允许规范提供详细的功能，元素通过\textit{std::move()}传递给可调用对象。在可调用对象内部，可以决定如何处理:\par

\begin{itemize}
  \item 按值接受参数，总是移动/窃取值或资源
  \item 通过右值/通用引用参数来决定移动/窃取哪个值/资源
\end{itemize}

例如:\par

\begin{codeboxt}[cpp]{lib/foreachmove.cpp}
#include <iostream>
#include <string>
#include <vector>
#include <algorithm>
template<typename T>
void print(const std::string& name, const T& coll)
{
  std::cout << name << " (" << coll.size() << " elems): ";
  for (const auto& elem : coll) {
    std::cout << " \"" << elem << "\"";
  }
  std::cout << "\n";
}
void process(std::string s) // gets moved value from rvalues
{
  std::cout << "- process(" << s << ")\n";
  ...
}
int main()
{
  std::vector<std::string> coll{"don't", "vote", "for", "liars"};
  print("coll", coll);
  // move away only the elements processed:
  std::for_each(std::make_move_iterator(coll.begin()),
  std::make_move_iterator(coll.end()),
  [] (auto&& elem) {
    if (elem.size() != 4) {
      process(std::move(elem));
    }
  });
  print("coll", coll);
}
\end{codeboxt}

代码中，将所有大小为4的元素移动到helper函数\textit{process()}中:\par

\begin{codebox}[cpp]
// move away only the elements processed:
std::for_each(std::make_move_iterator(coll.begin()),
std::make_move_iterator(coll.end()),
[] (auto&& elem) {
  if (elem.size() != 4) {
    process(std::move(elem));
  }
});
\end{codebox}

为此，将标记的元素(移动迭代器标记为\textit{std::move()})通过通用引用传递给\textit{process()}，并将其与\textit{std::move()}一起传递给\textit{process()}。因为\textit{process()}按值接受参数，所以值实际上已经移走。\par

因此，所有大小不为4的元素都从容器\textit{coll}中的对象中移出。该程序的输出如下(?表示未知值):\par

\begin{myOutput}
coll (4 elems): "don't" "vote" "for" "liars" \\
- process(don't) \\
- process(for) \\
- process(liars) \\
coll (4 elems): "?" "vote" "?" "?"
\end{myOutput}

最后一行是这样的:\par

\begin{myOutput}
coll (4 elems): "" "vote" "" ""
\end{myOutput}

这里使用了辅助函数\textit{std::\allowbreak{}make\_\allowbreak{}move\_\allowbreak{}iterator(\allowbreak{})\allowbreak{}}，这样在声明迭代器时就不必指定元素类型了。从C++17开始，类模板实参推导(CTAD)允许直接声明类型\textit{std::\allowbreak{}move\_\allowbreak{}iterator}，而不需要指定元素类型:\par

\begin{codebox}[cpp]
std::for_each(std::move_iterator{coll.begin()},
std::move_iterator{coll.end()},
[] (auto&& elem) {
  if (elem.size() != 4) {
    process(std::move(elem));
  }
});
\end{codebox}

\mySmallsection{14.3.2 在构造函数和成员函数中移动迭代器}

也可以在只读取一次元素的算法中使用移动迭代器，场景可能是将源容器的元素移动到另一个容器(相同或不同类型的)。例如:\par

\begin{codeboxt}[cpp]{lib/moveitor.cpp}
#include <iostream>
#include <string>
#include <list>
#include <vector>

template<typename T>
void print(const std::string& name, const T& coll)
{
  std::cout << name << " (" << coll.size() << " elems): ";
  for (const auto& elem : coll) {
    std::cout << " \"" << elem << "\"";
  }
  std::cout << "\n";
}

int main()
{
  std::list<std::string> src{"don't", "vote", "for", "liars"};
  // move all elements from the list to the vector:
  std::vector<std::string> vec{std::make_move_iterator(src.begin()),
    std::make_move_iterator(src.end())};
  print("src", src);
  print("vec", vec);
}
\end{codeboxt}

该程序有以下输出:(?表示未知值):\par

\begin{myOutput}
src (4 elems): "?" "?" "?" "?" \\
vec (4 elems): "don't" "vote" "for" "liars"
\end{myOutput}

请再次注意，源容器中的元素数量没有改变。所有元素会移动到初始化的新容器中。因此，源范围中的元素随后处于已移动状态，所以不知道它们的值。\par


\section{移动迭代器}

通过使用移动迭代器(在C++11中引入)，甚至可以在其他算法中使用移动语义，也可以在任何输入范围中使用移动语义(例如，在构造函数中)。\par

但是，使用这些迭代器时要小心。在迭代容器或范围的元素时，每次对元素的访问都使用\textit{std::move()}。这可能会更快，会使元素处于有效但未定义的状态，所以一个元素不应该使用两次\textit{std::move()}。\par


\section{总结}

\begin{itemize}
  \item C++标准库提供了移动多个元素的特殊算法，称为\textit{std::move()}(是\textit{std::move()}的重载，用于将单个对象标记为可移动)和\textit{std::\allowbreak{}move\_\allowbreak{}backward(\allowbreak{})\allowbreak{}}。它们让元素处于离开状态。
  \item 移除算法可以将元素保留在已移动状态。
  \item 移动迭代器允许在迭代元素时使用移动语义。可以在算法或构造函数中使用这些迭代器，构造函数中，范围用于初始化/设置值。但是，要确保迭代器只对每个元素使用一次。
\end{itemize}
